\documentclass[11pt]{article}
\usepackage[a4paper,margin=21mm]{geometry}
\usepackage[no-math]{fontspec}
\setmainfont{Latin Modern Roman}
\newfontfamily\CJKfont{Noto Serif CJK SC}
\XeTeXlinebreaklocale "zh"
\XeTeXlinebreakskip=0pt plus 1pt
\usepackage{amsmath,amssymb,amsthm,amscd,graphicx,color}
\usepackage{xurl}
\usepackage[unicode,hidelinks]{hyperref}
\allowdisplaybreaks
\setlength\emergencystretch{3em}
\numberwithin{equation}{section}

\numberwithin{equation}{section}

\newtheorem{Theorem}{Theorem}[section]
\newtheorem*{Theorem*}{Theorem}
\newtheorem{Corollary}[Theorem]{Corollary}
\newtheorem{Lemma}[Theorem]{Lemma}
\newtheorem{Proposition}[Theorem]{Proposition}
 { \theoremstyle{definition}
\newtheorem{Definition}[Theorem]{Definition}
\newtheorem{Note}[Theorem]{Note}
\newtheorem{Example}[Theorem]{Example}
\newtheorem{Remark}[Theorem]{Remark} }

\newcommand{\C}{{\mathbb C}}
\DeclareMathOperator{\End}{End}
\DeclareMathOperator{\dimension}{dim}
\DeclareMathOperator{\tr}{tr}
\DeclareMathOperator{\id}{id}
\DeclareMathOperator{\inv}{inv}
\newcommand{\Uqgl}{{U_q(\mathfrak{gl}(N+1))}}
\newcommand{\ee}{{ \hat{e}}}
\newcommand{\E}{{ \hat{E}}}
\newcommand{\jo}{{(j_0;\dots;j_v;\dots;j_N)}}
\newcommand{\zz}{{\mu}}
\newcommand{\ii}{{\mathbf{i}}}
\newcommand{\jj}{{\mathbf{j}}}
\newcommand{\mmm}{{\mapsto}}
\newcommand{\zo}{{(\mu)}}
\newcommand{\pp}{{P_2^+(V^{(2)})}}

\begin{document}
\begin{center}\Large Explicit fused symmetric-power central elements of generic Uq(gl(N+1)) in modified root-vector generators\end{center}
\noindent\textbf{Stable ID:} 1093\quad\textbf{Author:} Jeffrey Kuan; Keke Zhang\par
\noindent\textbf{Work:} Explicit Central Elements of $U_q(\mathfrak{gl}(N+1))$\par
\noindent\textbf{Source:} \url{https://sigma-journal.com/2023/036/}\par
\noindent\textbf{Version:} Official SIGMA 19 (2023) article 036, DOI 10.3842/SIGMA.2023.036; journal PDF and native TeX acquired 2026-09-30, exact bytes pinned by SHA256\par
\noindent\textbf{Rights:} CC BY 4.0. Authors retain copyright. \url{https://creativecommons.org/licenses/by/4.0/}. Reuse and typesetting adaptation; original language and mathematical content preserved.\par
\noindent\textbf{Difficulty (prior selection preserved):} {\CJKfont H2: The new proof evaluates fused R-matrices explicitly on a q-symmetrized basis rather than leaving the central elements as an abstract quantum trace. It proves invariance of the root-vector sums under minimal coset representatives, constructs the symmetric-power basis by inversion induction, and factors the fused operator into the stated root-vector coefficients. These combinatorial and representation calculations are the substantive obligation before the established quantum-trace centrality construction can be used.}\par
\medskip\noindent\textbf{Editorial boundary note (not author text):} Complete original main Theorem3.1, with the entire current explicit fused-operator evaluation. All symmetric-group/coset/reversal/weight conventions and current coset connectivity/weight-identity proofs, full quantum-group relations, root-vector definitions and stated commutation identities, symmetric-power/fused R-matrix and quantum-trace conventions, both coefficient-sum invariance statements and full current proof, complete q-symmetric basis construction by inversion induction, matrix-unit identities and full fused factorization/trace computation are retained. Original Remark3.4 and the explicit normalized representation are included to preserve the original root-relation footnote and coefficient conventions. Both counterpart computations that the author calls similar remain with their actual displayed identities and supplied inputs. Jimbo fusion and Drinfeld quantum-trace centrality are precise author-stated formal prerequisites, not new constructions replaced by editor proof. This is one main explicit-central-element outcome; supporting basis/coset statements and optional examples/character formulas are not separate counts. Literal author primes, reversed products, signs and index spelling are retained.\par
\medskip\hrule\medskip\noindent\textbf{Author text begins below.}\par
\setcounter{section}{1}
% Original source lines 76--83
The explicit generators of the center of $\Uqgl$ was first constructed for generic $q$
in~\cite{FaddeevRT} without proof (see~\cite{Arnaudon:1993yg}). Previous work of~\cite{ZGBCMP} applies Drinfeld's central element construction~\cite{Drin90} to universal $R$-matrices in order to construct central elements of quantum groups and to determine their eigenvalues on irreducible highest weight modules.

Jimbo in \cite{Jimbo1986} provided a method to construct a polynomial $z(x) = \sum_{i=0}^{N+1} z_ix^i$ in variable~$x$ such that all $z_i$’s are central. Tanisaki in \cite{Tanisaki1992KILLINGFH} proved that all these $z_i$’s generate the whole centre by quantized Harish-Chandra isomorphism. However, the expression of $z_i$’s in \cite{Jimbo1986} are not explicit enough for applications cited in the first paragraph.

By using Jimbo's \cite{Jimbo1986} formula for the $R$-matrix of $\Uqgl$, further work by the same authors~\cite{GZB} explicitly writes a quantum Casimir element of $\Uqgl$ with a formula for its eigenvalues. Using an explicit formula for the universal $R$-matrix in~\cite{KirillovResh}, the authors~\cite{ZGBJPhysA} write an explicit (but somewhat complicated) expression for general Casimir elements in quantum groups.

In this paper, we apply Drinfeld's central element construction to the fused $R$-matrices of $\Uqgl$ in~\cite{Jimbo1986}, rather than the universal $R$-matrices of~\cite{KirillovResh}. The resulting central elements appear to be slightly simpler than the previous expressions. The proof requires some new ingredients, notably relations between the root vectors~\cite{KT91,Lusztig90,Xi1994} and some elementary knowledge of coset representatives of symmetric groups.

\setcounter{section}{1}
% Original source lines 87--332
\section{Notations and backgrounds}
\subsection{Symmetric groups}
Recall that the symmetric group is the Weyl group of the special linear group. Define the usual action of the symmetric group $S_m$ on $\mathbb{N}^m$ by $ \sigma(x_1,\dots,x_m)= (x_{\sigma(1)}, \dots, x_{\sigma(m)})$. For any $A\in \mathbb{N}^m$, define $ H_A \le S_m $ to be the subgroup $ \{\sigma \in S_m\colon \sigma(A)=A\}$. Let $ D_A \subset S_m $ be the set of (left) coset representatives of $H_A$ with the fewest inversions: in other words, $ \sigma \in D_A $ if and only if $ \inv (\sigma) \le \inv (\tau) $ for every $ \tau \in \sigma H_A$. Here, $\inv$ is the length function defined as the fewest number
of simple reflections constituting the Weyl group element.

 \begin{Example} Take $ A= (3,3,3,2,2,1).$ Then $H_A=S_3 \times S_2 \times S_1 \leq S_6$. We have $ (34) \in D_A $ with $ (34) \cdot A= (3,3,2,3,2,1)$, and $ (35) \in D_A $ with $ (35)\cdot A=(3,3,2,2,3,1) $. However, $ (134)\cdot A= (34) \cdot A$ and $\inv((134)) > \inv((34))$, so $ (134) \notin D_A $.
 \end{Example}

 We recall (see, e.g., \cite{Carter}) that each coset of $H_A$ has a unique representative $\sigma \in D_A$, and that $\inv(\sigma\tau) = \inv(\sigma) + \inv(\tau)$ for every $\tau \in H_A$.

 \begin{Lemma}\label{Gene}
 	Suppose that $\tau,\tau' \in D_A$. Then there exist a sequence of elements $\tau_0,\dots,\tau_l\in D_A$ such that $\tau_0 = \tau'$, $\tau_l=\tau$ and every $\tau_{j+1}\tau_j^{-1}$ is a transposition for $j=0,\dots,l-1$.
 \end{Lemma}
 \begin{proof}
 	It suffices to prove this statement when either $\tau$ or $\tau'$ is the identity permutation $e$, because the two sequences can be concatenated. So suppose that $\tau$ is an arbitrary element of $D_A$ and $\tau'=e$. Let $s_k\dots s_1$ be a minimal word representation of $\tau$ and set $\tau_j = s_j \dots s_1$. Then $\tau_{j+1}\tau_j^{-1}=s_{j+1}$, which is a transposition. So it remains to show that $\tau_j \in D_A$. If it were not, then there would exist a transposition $s\in H_A$ such that $\inv(\tau_js) = \inv(\tau_j)-1$. But then $\inv(\tau s) = \inv(\tau)-1$, contradicting the assumption that $\tau \in D_A$.
 	
 	Now suppose that $\tau=e$ and $\tau'$ is an arbitrary element of $D_A$. By the previous paragraph, there exist a sequence of elements $\tilde{\tau}_0=e,\tilde{\tau}_1,\dots,\tilde{\tau}_l=\tau'$ in $D_A$ such that every $\tilde{\tau}_{j+1}\tilde{\tau}_j^{-1}$ is a~transposition. Setting $\tau_j = \tilde{\tau}_{l-j}$, we have that $\tau_0=\tau',\dots,\tau_l=e$ is a sequence of elements in $D_A$ and ${\tau}_{l-j-1}{\tau}_{l-j}^{-1}$ is a transposition for every $j$. The latter equality is equivalent to the condition that for every $k$, $\tau_k = s\tau_{k+1}$ for some transposition $s$. Since this is also equivalent to the condition that $\tau_{k+1}\tau_k^{-1}$ is a transposition for every $k$, this finishes the proof.
 \end{proof}

 Let $ \mathcal{B}_m^{(N)} $ denote the set of sequences $ \mu=(\mu_0,\dots,\mu_N) $ of non-negative integers such that $ \mu_0+\dots+\mu_N=m $. For any $ \mu \in \mathcal{B}_m^{(N)} $, let $ H^\mu \le S_m$ denote the subgroup $ S_{\mu_0}\times \dots \times S_{\mu_N} $, and likewise let $ D^\mu $ denote the set of left coset representatives with the fewest inversions. Define $\mathcal{B}_m$ to be the union $\bigcup_{N\geq 1} \mathcal{B}_m^{(N)}$.

Let $\mathcal{W}_m \subset \mathbb{N}^m$ denote the subset of elements $(i_1,\dots,i_m)$ satisfying $i_1 \leq \dots \leq i_m$. For $\ii \in \mathcal{W}_m$, and assuming that $N\geq i_m$, let $\mu^{(N)}(\ii) \in \mathcal{B}_m^{(N)}$ be defined by
\[
\mu^{(N)}(\ii)_k = \vert \{ l \in \{1,\dots,m\}\colon i_l=k \} \vert.
\]
For $\ii \in \mathcal{W}_m$ satisfying $i_m \leq N$, we have a natural isomorphism between the subgroups $H_{\ii}$ and $H^{\mu^{(N)}(\ii)}$. Thus there is also a natural bijection between $D_{\ii}$ and $D^{\mu^{(N)}(\ii)}$.

Given $\ii \in \mathcal{W}_m$, define the equivalence relation $\sim_{\ii}$ on $S_m$ by
\[
\tau \sim_{\ii} \sigma \text{ if and only if } \sigma^{-1}\tau \in H_{\ii}.
\]
In words, $\tau \sim_{\ii} \sigma$ if and only if they are in the same left coset of $H_{\ii}$. For any $\tau \in S_m$, there is a unique $\sigma \in D_{\ii}$ such that $\tau \sim_{\ii} \sigma$. Define $d_{\ii}(\tau)= \inv\big(\sigma^{-1}\tau\big)$. In other words, if $\tau$ is written uniquely as $\tau=\sigma \xi$ for $\sigma \in D_{\ii}$ and $\xi \in H_{\ii}$, then $d_{\ii}(\tau) = \inv(\xi)$. We will also let $\sigma_{\ii}(\tau)$ and $\xi_{\ii}(\tau)$ denote the two permutations in the unique expression $\tau = \sigma \xi$.

Finally, given $\tau \in S_m$, let $\bar{\tau}$ denote the reversed permutation $\bar{\tau}(k)= \tau(m+1-k)$.

We conclude this section by noting the following identity.
\begin{Lemma}\label{Leem}
For $\ii \in \mathcal{W}_m$, set $\mu = \mu^{(N)}(\ii)$. Then
\[
-N\mu_{0} -(N-2)\mu_{1} + \dots + N\mu_{N} = -Nm + 2(i_1+ \dots + i_m).
\]
\end{Lemma}
\begin{proof}
By definition,
\[
\mu_k = | l \in \{1,\dots,m\}\colon i_l=k |.
\]
We thus rewrite
\begin{align*}
-N\mu_{0} -(N-2)\mu_{1} + \dots + N\mu_{N} &= -N(\mu_0 + \dots + \mu_N) + 2\mu_1 + 4\mu_2 + \dots + 2N\mu_N \\
&= -Nm+ 2\mu_1 + 4\mu_2 + \dots + 2N\mu_N
\end{align*}
and note that
\[
2(i_1 + \dots + i_m) = 2\sum_{k=1}^N k\cdot | l \in \{1,\dots,m\}\colon i_l=k | = 2\sum_{k=1}^N k\cdot \mu_k,
\]
which shows the identity.
\end{proof}

\subsection{Quantum groups}
We use the following notation modified from \cite{Jimbo1986}.
Define $ U_q(\mathfrak{sl}(N+1)) $ to be the associative algebra over $ \mathbb{C} $ generated by the symbols $ q^{\pm h_i/2}$, $\hat{e}_{\pm, i}$, $1 \leq i\leq N$, under the following relations:
\begin{gather*}
q^{h_i/2} \cdot q^{-h_i/2}= q^{-h_i/2} \cdot q^{h_i/2}=1, \qquad q^{h_i/2} \cdot q^{h_{i'}/2}= q^{h_{i'}/2} \cdot q^{h_i/2},\nonumber\\
 q^{h_i/2} \hat{e}_{\pm, i'} q^{-h_i/2}= q^{\pm a_{ii'}/2} \hat{e}_{\pm, i'},\nonumber\\
[\hat{e}_{+,i}, \hat{e}_{-,i'}] =\delta_{i,i'} \frac{q^{h_i}-q^{-h_i}}{q-q^{-1}}, \nonumber\\
 \hat{e}_{\pm, i}\hat{e}_{\pm, i'} = \hat{e}_{\pm ,i'}\hat{e}_{\pm ,i}, \qquad |i-i'|\geq 2, \nonumber\\
 \hat{e}_{\pm, i}^2 \hat{e}_{\pm, i\pm 1} -\big(q+q^{-1}\big) \hat{e}_{\pm, i} \hat{e}_{\pm ,i\pm 1} \hat{e}_{\pm ,i} + \hat{e}_{\pm, i\pm 1} \hat{e}_{\pm, i}^2=0,\qquad 1\leq i, i\pm 1 \leq N. %\label{UqglDef}
\end{gather*}
Here, $ (a_{i,i'})_{1\leq i,i' \leq N} $ denotes the Cartan matrix of type $ A_N $, i.e.,
\begin{equation*}
a_{ii'}=
\begin{cases}
\hphantom{-}2,& i = i',\\ -1, & i=i' \pm 1,\\ \hphantom{-}0,& \text{otherwise}.
\end{cases}
\end{equation*}
Then define $ U_q(\mathfrak{gl}(N+1)) $ by adjoining to $ U_q(\mathfrak{sl}(N+1)) $ the elements $ q^{\pm \epsilon_i/2}$, $0 \leq i \leq N $, so that $ q^{h_i}= q^{\epsilon_{i-1}-\epsilon_i} $ and that $ q^{\epsilon_0+\dots +\epsilon_N} $ belongs to the center.

The $m$-fold co-product is the algebra homomorphism \[\Delta^{(m)}\colon \ \Uqgl \rightarrow \underbrace{\Uqgl \otimes \dots \otimes \Uqgl}_{m} \]
such that \begin{gather}
 \Delta^{(m)}\big (q^{\pm \epsilon_i/2}\big)= q^{\pm \epsilon_i/2} \otimes \dots \otimes q^{\pm \epsilon_i/2},\qquad
 \Delta^{(m)} (\hat{e}_{\pm ,i} ) =\sum_{v=1}^{m} \hat{e}^{(v)}_{\pm ,i} ,\label{qerel}
\end{gather}
where
\[
\hat{e}^{(v)}_{\pm ,i} = \underbrace{q^{h_i/2} \otimes \dots \otimes q^{h_i/2}}_{v-1} \otimes \hat{e}_{\pm ,i} \otimes \underbrace{ q^{-h_i/2} \otimes \dots \otimes q^{-h_i/2}}_{m-v}.
\]
 We also have the reversed co-product
 \begin{gather*}
 \bar{\Delta}^{(m)} \big(q^{\pm \epsilon_i/2}\big)= q^{\pm \epsilon_i/2} \otimes \dots \otimes q^{\pm \epsilon_i/2}, \\
 \bar{\Delta}^{(m)} (\hat{e}_{\pm, i})= \sum_{v=1}^{m} \underbrace{q^{-h_i/2} \otimes \dots \otimes q^{-h_i/2}}_{v-1} \otimes \hat{e}_{\pm, i} \otimes q^{h_i/2} \otimes \dots \otimes q^{h_i/2}.
 \end{gather*}
We write $\Delta$, $\bar{\Delta}$ for $\Delta^{(2)}$, $\bar{\Delta}^{(2)}$. The map $ \Delta $ endows $ \Uqgl $ with a structure of a bi-algebra. It is also a Hopf algebra, but we will not need the counit and antipode.

Consider the following relations for $R$:
\begin{gather}
 R \Delta(u) = \bar{\Delta}(u) R \qquad \forall u\in \Uqgl. \label{Equi}
\end{gather}
 This admits a unique (up to a multiplicative constant) solution
$ R \in \End (V\otimes V) $, where $ V:=\mathbb{C}^{N+1} $ is the defining representation of $ \Uqgl $.

 Let us consider \eqref{Equi} in $ \Uqgl \otimes \End \big(\C^{N+1}\big) $. Then a family of solutions is given by (see \cite{Jimbo1986}) $R(x)= \sum_{0 \leq i,j \leq N} \hat{E}_{ij}'(x) \otimes e_{ji} $, where
 \begin{gather*} %\label{ehatx}
 \hat{E}'_{ij}(x) = \begin{cases}
 \big(x q^{(\epsilon_i+\epsilon_j-1)/2}\big)^{\mp 1} {E}'_{ij} , & i \lessgtr j,\\
 \big(xq^{\epsilon_i}- x^{-1}q^{- \epsilon_i}\big)/\big(q-q^{-1}\big), & i=j.
 \end{cases}
 \end{gather*}
 For this paper, only the value at $x=1$ will be used in the proofs, but for completeness Jimbo's result is stated in its full generality. Here, $ E'_{ij} $ are the root vectors defined recursively by
 \begin{gather*}
E'_{ij}= E_{ik}'E_{kj}' -q^{\pm 1} E_{kj}'E_{ik}',\qquad i \lessgtr k \lessgtr j, \qquad E'_{i-1,i} = \hat{e}_{+,i}, \qquad E'_{i,i-1} = \hat{e}_{-,i}
\end{gather*}
and $e_{ji}$ is the usual matrix which acts on the canonical basis $\{I_0,\dots,I_N\}$ of $\mathbb{C}^{N+1}$ by
\[
e_{ji}I_k = \delta_{ik} I_j.
\]

In \cite{GZB}, the authors define
\[
\hat{E}_{ij}
=
\begin{cases}
q^{(E_{ii}+E_{jj}-1)/2}E_{ij}, & i \neq j,\\
q^{E_{ii}}/\big(q-q^{-1}\big), & i=j,
\end{cases}
\]
where the modified root vectors are
\begin{equation*}
E_{ij}= E_{ik}E_{kj} -q^{- 1} E_{kj}E_{ik},\qquad i \lessgtr k \lessgtr j, \qquad E'_{i-1,i} = \hat{e}_{+,i}, \qquad E'_{i,i-1} = \hat{e}_{-,i}.
\end{equation*}
Using Jimbo's results, they show that
\begin{gather*}
R = \sum_{N \geq i\geq j \geq 0} \hat{E}_{ij} \otimes e_{ji} \in \Uqgl \otimes \End (V),\nonumber\\
R^T = \sum_{N \geq i\geq j \geq 0} \hat{E}_{ji} \otimes e_{ij} \in \Uqgl \otimes \End (V) %\label{r}
\end{gather*}
satisfy
\begin{gather*}
R\Delta(u) = \bar{\Delta}(u) R,\qquad
R^T\bar{\Delta}(u) = {\Delta}(u) R^T.
\end{gather*}

We can solve \eqref{Equi} even more generally. Consider the $ m $-fold tensor of the defining representation $ V^{\otimes m}$, and let $P_mV^{\otimes m}$ be the symmetric projection. In other words, $P_mV^{\otimes m}$ is isomorphic to the module of highest weight~$m\nu$, where~$\nu$ is the highest weight of~$V$.

Then the solution of \eqref{Equi} in $ \Uqgl \otimes \End (P_m V^{\otimes m}) $ is given by the fused $R$-matrix~\cite{Jimbo1986}%
\begin{gather*} %\label{fusx}
R^{}_{0,\{1,2,\dots,m\}}(x)= R_{0m}(x)R_{0m-1}(xq) \cdots R_{01}\big(xq^{ m-1}\big) P_m \in \Uqgl\otimes \End \big(P_m V^{\otimes m}\big) .
\end{gather*}

Therefore, the fused $R$-matrices
\begin{gather*}
R^{}_{0,\{1,2,\dots,m\}} = R_{0m}R_{0m-1}\cdots R_{01}P_m \in U_q(\mathfrak{gl}_{N+1})\otimes \End \big(P_m V^{\otimes m}\big), \\
R^{T}_{0,\{1,2,\dots,m\}} = R^T_{0m}R^T_{0m-1}\cdots R^T_{01}P_m \in U_q(\mathfrak{gl}_{N+1})\otimes \End \big(P_m V^{\otimes m}\big)
\end{gather*}
satisfy
\begin{gather}
R_{0,\{1,2,\dots,m\}} \Delta(u) = \bar{\Delta}(u) R_{0,\{1,2,\dots,m\} },\qquad
R^T_{0,\{1,2,\dots,m\}} \bar{\Delta}(u) = {\Delta}(u) R^T_{0,\{1,2,\dots,m\}}\label{RRel}
\end{gather}
in $\Uqgl \otimes \End (P_m V^{\otimes m})$ for all $u\in \Uqgl$.

Therefore,
\begin{equation*}
\Gamma_m: = R^T_{0,\{1,2,\dots,m\}} R_{0,\{1,2,\dots,m\}} \in \Uqgl \otimes \End \big(P_m V^{\otimes m}\big)
\end{equation*}
commutes with $ \Delta (u) $ for all $u\in \Uqgl$. By Drinfeld's central element construction \cite{Drin90}, the element
\begin{equation*}
C_m= \id \otimes\tr_q(\Gamma_m)
\end{equation*}
is central in $\Uqgl$, where the quantum trace $ \tr_q $ of an operator $A$ is defined by
\begin{equation*}%\label{key}
\tr_q(A) = \tr\big(q^{-2h_{\rho}}A\big) := \tr \big(q^{-N\epsilon_0-(N-2)\epsilon_1 - \dots +N\epsilon_N}A\big).
\end{equation*}

We also have the following relations between the root vectors:\footnote{Similar relations can be found in the paper \cite{Xi1994}, but there appear to be some typos. They can also be derived from \eqref{RRel} by applying $\id \otimes B$ to both sides, for suitable linear maps $B$ on $\mathrm{End}\big(P_2V^{\otimes 2}\big)$. One can also check that the relations hold in the explicit representations in Remark~\ref{Remark1} below.}
For $i<l$ and $j<k$,
\begin{gather}\label{comm}
{E}_{il}{E}_{jk} =
\begin{cases}
{E}_{jk}{E}_{il}, & i< j < k < l \text{ or } i<l<j<k,\\
q^{-1} {E}_{jk}{E}_{il}, & i=j<k<l, \\
q {E}_{jk}{E}_{il}, & i<j<k=l,\\
{E}_{jk}{E}_{il} + \big(q-q^{-1}\big) {E}_{jl}{E}_{ik}, & i<j<l<k,
\end{cases}
\\
\label{comm1}
{E}_{li}{E}_{kj} =
\begin{cases}
{E}_{kj}{E}_{li}, & i< j < k < l \text{ or } i<l<j<k,\\
q^{} {E}_{kj}{E}_{li}, & i=j<k<l, \\
q^{-1} {E}_{kj}{E}_{li}, & i<j<k=l,\\
{E}_{kj}{E}_{li} - \big(q-q^{-1}\big) {E}_{lj}{E}_{ki}, & i<j<l<k.
\end{cases}
\end{gather}
By the first and fourth lines above,
\begin{gather}\label{comm2}
\hat{E}_{il}\hat{E}_{jk} + q^{-1} \hat{E}_{ik}\hat{E}_{jl} = \hat{E}_{jk}\hat{E}_{il} + q \hat{E}_{jl}\hat{E}_{ik}, \\
\hat{E}_{li}\hat{E}_{kj} + q^{} \hat{E}_{ki}\hat{E}_{lj} = \hat{E}_{kj}\hat{E}_{li} + q^{-1} \hat{E}_{lj}\hat{E}_{ki}. \label{comm3}
\end{gather}

For any $ \ii $ and $ \jj $ in $ \mathbb{N}^m $, define
\[
\hat{E}^{\pm}_{\ii \jj} :=
\begin{cases}
\hat{E}_{i_1j_1}\cdots \hat{E}_{i_mj_m}, & \pm(i_1-j_1),\dots,\pm(i_m-j_m)\geq 0,\\
0, & \text{else}.
\end{cases}
\]
 and let $ e_{ \jj \ii }:= e_{j_1i_1} \otimes \dots \otimes e_{j_mi_m}$. For $\ii, \jj \in \mathcal{W}_m$, define
\[
\tilde{E}^{\pm}_{\ii\jj} =
\sum_{\zeta \in D_{\ii}}q^{\pm(\inv(\tau) -\inv(\zeta))}\hat{E}^{\pm}_{\bar{\zeta}(\ii)\bar{\tau}(\jj)},
\]
where $\tau$ is an arbitrary element of $D_{\jj}$. Note that the notation does not depend on $\tau$, which is justified by the following lemma and the relation $\inv(\bar{\tau}) = (m-1)m/2 - \inv(\tau)$.

\begin{Lemma} \label{prop 1}\quad
	\begin{enumerate}\itemsep=0pt
		\item[$1.$] For every $\tau,\tau' \in D_A$, the following identity holds: \begin{equation*}%\label{key}
		\sum_{\sigma\in D_B} q^{\mp (\inv(\sigma)-\inv(\tau))} \hat{E}^{\pm}_{\tau(A) \sigma (B)} = \sum_{\sigma\in D_B} q^{\mp (\inv(\sigma)-\inv(\tau'))} \hat{E}^{\pm}_{\tau'(A) \sigma (B)} .
		\end{equation*} 	 	
		\item[$2.$] Likewise, for every $\sigma,\sigma' \in D_B$, \begin{equation*}%\label{key}
		\sum_{\tau\in D_A} q^{\mp (\inv(\tau)-\inv(\sigma))} \hat{E}^{\pm}_{\tau(A) \sigma (B)} = \sum_{\tau\in D_A} q^{\mp (\inv(\tau)-\inv(\sigma'))} \hat{E}^{\pm}_{\tau(A) \sigma' (B)} .
		\end{equation*} 	 	
	\end{enumerate}
	\end{Lemma}
\begin{proof}
	For $m=2$, both cases are equivalent to the relations in \eqref{comm} and \eqref{comm1}.
	
	Now suppose that $ m>2$. We only prove part 1, as part 2 is similar. By Lemma \ref{Gene}, it suffices to consider the case when $\tau'=s\tau$ where $s$ is a transposition, and assume without loss of generality that $\inv(\tau')= \inv(\tau)+1$. Define the two sets $D_B^{(1)}$ and $D_B^{(2)}$ by
	\begin{gather*}
		D_B^{(1)}= \{ \sigma \in D_B \colon s \sigma(B)=B \},\qquad
		D_B^{(2)}= \{ \sigma \in D_B \colon s \sigma(B) \neq B \}.
	\end{gather*}	
	First, note that by the second and third lines of \eqref{comm},
\[
	\sum_{\sigma\in D_B^{(1)}} q^{\inv(\sigma)-\inv(\tau)} \hat{E}^-_{\tau(A) \sigma (B)} = \sum_{\sigma\in D_B^{(1)}} q^{\inv(\sigma)-\inv(s\tau)} \hat{E}^-_{s \tau(A) \sigma (B)} .
\]
	Partition the set $D_B^{(2)}$ into two sets $D_-$ and $D_+$ of equal cardinality, where $\sigma \in D_-$ if and only if $\inv(\sigma) < \inv(s \sigma)$. Then
	\begin{gather*}
	\sum_{\sigma\in D_B^{(2)}} q^{\inv(\sigma)-\inv(\tau)} \hat{E}^-_{\tau(A) \sigma (B)} = \sum_{\sigma \in D_-} \big( q^{\inv(\sigma)-\inv(\tau)} \hat{E}^-_{\tau(A) \sigma (B)} + q^{\inv(\sigma)+1-\inv(\tau)} \hat{E}^-_{\tau(A) s\sigma (B)} \big) \\
\qquad \stackrel{\eqref{comm2}}{=} \sum_{\sigma \in D_-} \big( q^{\inv(\sigma)+1-\inv(\tau')} \hat{E}^-_{s\tau(A) s\sigma (B)} + q^{\inv(\sigma)-1-\inv(\tau)} \hat{E}^-_{s\tau(A) \sigma (B)} \big)\\
\qquad =\sum_{\sigma \in D_-} \big( q^{\inv(\sigma)+1-\inv(\tau')} \hat{E}^-_{\tau'(A) s\sigma (B)} + q^{\inv(\sigma)-\inv(\tau')} \hat{E}^-_{\tau'(A) \sigma (B)}\big)\\
\qquad =\sum_{\sigma\in D_B^{(2)}} q^{\inv(\sigma)-\inv(\tau')} \hat{E}^-_{\tau'(A) \sigma (B)} ,
	\end{gather*}
	as needed.
	The proof for $E^+$ is similar, where one uses \eqref{comm3} instead of \eqref{comm2}.
\end{proof}


% Original source lines 346--355
\section{Statements and proofs}

The main theorem is the following expression for central elements of $\Uqgl$.
\begin{Theorem}\label{Main}
	The element given by
\begin{equation*}
C_m=\sum_{\ii,\jj \in \mathcal{W}_m} q^{ 2i_1 + \dots + 2i_m-Nm} \tilde{E}_{\jj \ii}^-\tilde{E}_{\ii\jj}^+
\end{equation*}
is central in $ \Uqgl $.
\end{Theorem}

\setcounter{Theorem}{3}
% Original source lines 379--396
\begin{Remark}\label{Remark1}
In \cite{ZGBJPhysA}, the central element $C^{\Lambda_0}$ is defined, where $\Lambda_0$ is the highest weight of a finite-dimensional irreducible module $V(\Lambda_0)$. The construction there is similar to the one here, with the major difference being the use of explicit universal $R$-matrices in place of fused $R$-matrices. Although it is not necessarily simple to check directly that $C_m$ equals (up to a~constant) $C^{\Lambda_0}$ for $V(\Lambda_0)=P_mV^{\otimes m}$, it is straightforward to check that their eigenvalues are the same.

If $V(\Lambda_0)$ has distinct weights $\lambda_1,\dots,\lambda_r$ with multiplicities $d_1,\dots,d_r$, then the eigenvalue of~$C^{\Lambda_0}$ on an irreducible module with highest weight $\Lambda$ is given by
\[
\sum_{k=1}^r d_k q^{2(\Lambda + \rho,\lambda_k)}.
\]
Here, $\rho$ is half the sum of the positive roots, and $(\cdot,\cdot)$ is the usual invariant bilinear form on $\mathfrak{h}^*$. Now take $\Lambda_0$ to be the highest weight of $P_mV^{\otimes m}$; then the distinct weights are elements of $\mathcal{B}^{(N)}_m$ with multiplicity $1$. Therefore the eigenvalue is
\[
\sum_{\mu \in \mathcal{B}^{(N)}_m} q^{(2\rho,\mu)}q^{2(\Lambda,\mu)}.
\]

If $C_m$ acts on the same module $V(\Lambda)$, then its eigenvalue can be found by evaluating on the lowest weight vector, because then only the diagonal terms $(\ii=\jj)$ have a nonzero contribution. So the eigenvalue is
\[
\big(q-q^{-1}\big)^{-2m}q^{-Nm}\sum_{\ii \in \mathcal{W}_m} q^{ 2i_1 + \dots + 2i_m} q^{(2\mu^{(N)}(\ii),\Lambda)}.
\]
By Lemma \ref{Leem}, this is $\big(q-q^{-1}\big)^{-2m}q^{-Nm}$ times the eigenvalue of $C^{\Lambda_0}$.
\end{Remark}

\setcounter{Theorem}{5}
% Original source lines 410--542
\subsection[Basis for P\_mV\^\{otimes m\}]{Basis for $\boldsymbol{P_mV^{\otimes m}}$}

Before proving Theorem \ref{Main}, we will write a basis for the symmetric projection $P_mV^{\otimes m}$.

Let $ I_0,\dots,I_N $ be the canonical basis of $ V $, and define the action on $ V $ by 	
\begin{gather}
\hat{e}_{+,i}I_j=\delta_{ij}I_{j-1},\qquad
\hat{e}_{-,i}I_{j-1}=\delta_{ij}I_j,\qquad
q^{\pm \epsilon_i/2}I_j=q^{\pm \delta_ij/2}I_j.\label{21}
\end{gather}
For $\mu \in \mathcal{B}_m^{(N)}$, define the vector
\begin{equation*}%\label{key}
v\zo =I_0^{\otimes \mu_0} \otimes \dots \otimes I_N^{\otimes \mu_N}
\end{equation*}
Define
\begin{equation*}
M(\zz) = \sum_{\sigma \in D^{\mu}}q^{-\inv (\sigma)} \sigma (v(\mu)).
\end{equation*}
Here, as before, $S_m $ acts on $ V^{\otimes m} $ by permuting the order.
By an abuse of notation, for $\ii \in \mathcal{W}_m$ we define
\[
v(\ii) = v\big(\mu^{(N)}(\ii)\big), \qquad M(\ii) = M\big(\mu^{(N)}(\ii)\big),
\]
where $N \geq i_m$. Note that these definitions do not depend on the value of $N$.
We briefly note the identity
\begin{equation}\label{Brief}
e^{\sigma(a)}_{\pm,i} ( \sigma(v(\mu)) ) = \sigma\big( e^{(a)}_{\pm,i} ( v(\mu) ) \big) \qquad \text{for all } \sigma \in S_m.
\end{equation}

We now show that the set $ \{M(\mu)\}_{\mu\in \mathcal{B}^{(N)}_m} $ gives a basis for $ P_{m}V^{\otimes m} $. A more general statement appeared in~\cite{KIMRN} with a more complicated proof, but the expression here is more convenient for calculations.
\begin{Theorem}
	The set $ \{M(\mu)\}_{\mu\in \mathcal{B}^{(N)}_m} $ gives a basis for $ P_{m}V^{\otimes m} $.
\end{Theorem}
\begin{proof}
Note that $ \big|\mathcal{B}^{(N)}_m\big|= \dimension P_mV^{\otimes m} $ and $ \{M(\mu)\}_{\mu\in \mathcal{B}^{(N)}_m} $ is a linearly independent set, so it suffices to show that $ M(\mu) \in P_mV^{\otimes m} $ for all $ \mu \in \mathcal{B}^{(N)}_m $. To show that every $ M(\mu) $ is in $ P_{m}V^{\otimes m} $, it suffices to show that \begin{equation} \label{hati}
\Delta^{(m)}(\hat{e}_{-,i}) (M(\mu))=q^{\frac{1}{2}(\mu_{i-1}+\mu_i-1)} \big(1+q^{-2}+q^{-4}+\dots+q^{-2\mu_i}\big) \cdot M\big(\mu-\hat{i}\big),
\end{equation}
where $ \hat{i}=(0,\dots,0,1,-1,0,\dots,0) $ with the $ -1 $ occurring in the $ i $th position. By~\eqref{qerel} and~\eqref{21}, $ \Delta^{(m)} (\hat{e}_{-,i})(M(\mu)) $ is in the span of $ \{\tau(v(\mu-\hat{i}))\}_{\tau\in S_m}$. Let $ A(\tau) $ be the coefficients in the expansion
\[
 \Delta^{(m)} (\hat{e}_{-,i})(M(\mu)) = \sum_{\tau \in D^{\mu-\hat{i}}} A(\tau)\tau \big(v\big(\mu-\hat{i}\big)\big).
\]
It suffices to show that $A(\tau)= q^{\frac{1}{2}(\mu_{i-1}+\mu_i-1)}\big(1+q^{-2}+q^{-4}+\dots+q^{-2\mu_i}\big) q^{-\inv(\tau)}$ for all $\tau \in D^{\mu-\hat{i}}$. In fact, we show something stronger: for all $\tau \in D^{\mu-\hat{i}}$, there exist elements $\sigma^{(0)},\dots, \sigma^{(\mu_i)}\in D^{\mu}$ such that for $0 \leq j \leq \mu_i$,
\[
\hat{e}_{-,i}^{(a_j)} \big( q^{-\inv(\sigma^{(j)})} \sigma^{(j)}(v(\mu)) \big) = q^{\frac{1}{2}(\mu_{i-1}+\mu_i-1)} q^{-2j} q^{-\inv(\tau)} \tau\big(v\big(\mu- \hat{i}\big)\big).
\]

We will proceed by induction on the value of $\inv(\tau)$, using Lemma \ref{Gene}. The base case is when~$\tau$ is the identity permutation. Then it is straightforward to check that
\[
\sigma^{(j)} = (\mu_{[0,i-1]} \ \ \mu_{[0,i-1]}+1 \ \ \dots \ \ \mu_{[0,i-1]}+j),
\]
where $\mu_{[0,i-1]}= \mu_0 + \dots + \mu_{i-1}$, satisfies the necessary conditions.

Now fix $\tau \in D^{\mu - \hat{i}}$, and suppose that the induction hypothesis holds for $\tau$. Suppose that $\tilde{\tau} \in D^{\mu - \hat{i}}$ satisfies $\tilde{\tau}=s\tau$ for some transposition $s$, and assume without loss of generality that $\inv(\tilde{\tau}) = \inv(\tau) + 1$. Define
\[
\tilde{\sigma}^{(j)}
=
\begin{cases}
s \sigma^{(j)}, & \text{if } s \sigma^{(j)} \in D^{\mu}, \\
\sigma^{(j)}, & \text{else. }
\end{cases}
\]
We now aim to prove that
\begin{equation}\label{Aim}
\hat{e}^{(s(a_j))}_{-,i}\big(q^{-\inv(\tilde{\sigma}^{(j)})} \tilde{\sigma}^{(j)}(v(\mu)) \big) = q^{\frac{1}{2}(\mu_{i-1}+\mu_i-1)} q^{-2j} q^{-\inv(\tilde{\tau})} \tilde{\tau}\big(v\big(\mu- \hat{i}\big)\big).
\end{equation}
If $\tilde{\sigma}^{(j)} = s\sigma^{(j)}$, then \eqref{Aim} follows from \eqref{Brief} and the induction hypothesis. Now assume that $\tilde{\sigma}^{(j)} = \sigma^{(j)}$. Then $s= (a_j-1 \ \ a_j)$ and
\[
\tilde{\sigma}^{(j)} ( v(\mu) ) = \underbrace{I_* \otimes \dots \otimes I_* }_{a_j-2} \otimes I_{i-1} \otimes I_{i-1} \otimes \underbrace{I_* \otimes \dots \otimes I_*}_{m-a_j}.
\]
Using that
\[
 P \big( \big(\hat{e}_{i,-} \otimes q^{-h_i/2}\big)(I_{i-1} \otimes I_{i-1}) \big) = q^{-1}\big(q^{h_i/2} \otimes \hat{e}_{i,-}\big)(I_{i-1} \otimes I_{i-1}) ,
\]
where $P(x\otimes y)=y\otimes x$ is the permutation operator, we have that
\[
\hat{e}^{(s(a_j))}_{-,i}\big(q^{-\inv(\tilde{\sigma}^{(j)})} \tilde{\sigma}^{(j)}(v(\mu)) \big) = q^{-1}\hat{e}^{(a_j)}_{-,i}\big(q^{-\inv({\sigma}^{(j)})} {\sigma}^{(j)}(v(\mu)) \big).
\]
And now \eqref{Aim} follows from the induction hypothesis.
\end{proof}
\begin{Remark}\label{ExRep}
	For all $ \mu \in \mathcal{B}^{(N)}_m $, let $ \tilde{M}(\mu)=c(\mu)M(\mu) $ where $ c(\mu) $ is defined inductively by
	\begin{equation*}%\label{key}
	c(m,0,\dots,0)=1, \qquad
	\frac{c(\mu-i)}{c(\mu)}=\frac{q^{\frac{\mu_{i-1}+\mu_i+1}{2}}\big(1-q^{-2\mu_i}\big)}{q^{\mu_i+1}-q^{-\mu_i-1}} ,
	\end{equation*} then equation \eqref{hati} is equivalent to
\begin{equation*}%\label{key}
	\Delta^{(m)} (\hat{e}_{-,i}) \tilde{M}(\mu) =\frac{q^{\mu_i+1}-q^{-\mu_i-1}}{q-q^{-1}} \tilde{M}\big(\mu-\hat{i}\big).
	\end{equation*} In fact, one can show that
	\begin{gather*}
		\Delta^{(m)} (\hat{e}_{+,i}) \tilde{M}(\mu) =\frac{q^{\mu_{i-1}+1}-q^{-\mu_{i-1}-1}}{q-q^{-1}} \tilde{M}\big(\mu+\hat{i}\big),\\
		\Delta^{(m)} (\hat{e}_{-,i}) \tilde{M}(\mu) =\frac{q^{\mu_i+1}-q^{-\mu_i-1}}{q-q^{-1}} \tilde{M}\big(\mu-\hat{i}\big),\\
	 \Delta^{(m)} \big(q^{\pm h_i/2}\big) \tilde{M}(\mu) =q^{\pm\frac{1}{2}(\mu_{i-1}-\mu_i)} \tilde{M}(\mu)
	\end{gather*} defines a representation on $ P_mV^{\otimes m} $. This is equivalent to the representation in \cite[equation~(3)]{KMMO} and \cite[Lemma~3.1]{KIMRN}.
\end{Remark}
\begin{Corollary}\label{Core} Let $\ii,\jj\in\mathcal{W}_m$. For any $\tau,\sigma \in D_{\ii}$ and any $\zeta \in D_{\jj}$,
\[
\big(q^{\inv(\tau)}e_{\zeta(\jj) \tau(\ii)} - q^{\inv(\sigma)}e_{\zeta(\jj) \sigma(\ii)} \big) \big|_{P_{ m}^+V^{\otimes m} } =0.
\]
	
	Furthermore,
\[
\bigg( \sum_{\tau \in D_{\jj}} q^{-d_{\ii}(\tau)}e_{\tau(\jj)\tau(\ii)} \bigg) M(\ii) = M(\jj) .
\]
\end{Corollary}
\begin{proof}
For any $\mu \in \mathcal{B}^{(N)}_m$ not equal to $\mu(\ii)$, it is straightforward that
\[
\big(q^{\inv(\tau)}e_{\zeta(\jj) \tau(\ii)} - q^{\inv(\sigma)}e_{\zeta(\jj) \sigma(\ii)} \big) M(\mu)=0-0=0.
\]
On the other hand,
\begin{gather*}
\begin{split}
&\big(q^{\inv(\tau)}e_{\zeta(\jj) \tau(\ii)} - q^{\inv(\sigma)}e_{\zeta(\jj) \sigma(\ii)} \big) M(\mu(\ii))\\
&\qquad = q^{\inv(\tau)}e_{\zeta(\jj) \tau(\ii)} \big(q^{-\inv(\tau)} \tau(v(\mu(\ii)) )\big) - q^{\inv(\sigma)}e_{\zeta(\jj) \sigma(\ii)} \big(q^{-\inv(\sigma)} \sigma(v(\sigma(\ii)) )\big) \\
&\qquad = \zeta(v(\mu(\jj))) - \zeta(v(\mu(\jj)))=0.
\end{split}
\end{gather*}

For the second statement, we use that
\[
e_{\tau(\jj)\tau(\ii)}\sigma(v(\ii))
=
\begin{cases}
\tau(v(\jj)), & \tau \sim_{\ii} \sigma, \\
0, & \text{else.}
\end{cases}
\]
in order to show
\begin{align*}
\bigg( \sum_{\tau \in D_{\jj}} q^{-d_{\ii}(\tau)}e_{\tau(\jj)\tau(\ii)} \bigg) \sum_{\sigma \in D_{\ii}} q^{-\inv(\sigma)} \sigma(v(\ii)) &= \sum_{\substack{\tau \in D_{\jj},\sigma \in D_{\ii} \\ \tau\sim_{\ii} \sigma}}q^{-d_{\ii}(\tau)+ \inv(\sigma)} e_{\tau(\jj)\tau(\ii)}\sigma(v(\ii))\\
&= \sum_{\tau \in D_{\jj}} q^{-\inv(\tau)} \tau(v(\jj)).\tag*{\qed}
\end{align*}\renewcommand{\qed}{}
\end{proof}


% Original source lines 553--602
With Corollary \ref{Core} as motivation, define for any $\ii,\jj \in \mathcal{W}_m$,
\[
\tilde{e}_{\jj \ii} = \sum_{\tau \in D_{\jj}} q^{-d_{\ii}(\tau)}e_{\tau(\jj)\tau(\ii)} .
\]

\subsection{Proof of Theorem \ref{Main}}
We will now prove Theorem \ref{Main}. Begin by rewriting $R_{0,\{1,\dots,m\}}$. By definition,
\begin{align*}
R_{0,\{1,\dots,m\}} &= R_{0m}R_{0m-1}\cdots R_{01} \\
&=\bigg( \sum_{i_m \geq j_m} \hat{E}_{i_mj_m} \otimes \id^{\otimes m-1} \otimes e_{j_mi_m} \bigg) \cdots \bigg(\sum_{i_1\geq j_1}\hat{E}_{i_1j_1}\otimes e_{j_1i_1} \otimes \id^{\otimes m-1}\bigg)
\\&= \sum_{i_m \geq j_m} \cdots \sum_{i_1\geq j_1} \hat{E}_{i_mj_m} \cdots \hat{E}_{i_1j_1} \otimes e_{j_1i_1} \otimes \dots \otimes e_{j_mi_m} \\
&= \sum_{\ii,\jj \in \mathcal{W}_m } \sum_{\zeta \in D_{\ii},\tau \in D_{\jj}}\hat{E}_{\bar{\zeta}(\ii)\bar{\tau}(\jj)}^{+} \otimes e_{\tau(\jj)\zeta(\ii)}.
\end{align*}
From here, the goal is to write this expression as an element of $\Uqgl \otimes \End\big(P_m\big(V^{\otimes m}\big)\big)$. By Corollary \ref{Core},
\[
R_{0,\{1,\dots,m\}} = \sum_{\ii,\jj \in \mathcal{W}_m } \sum_{\zeta \in D_{\ii},\tau \in D_{\jj}}q^{\inv(\sigma_{\ii}(\tau)) -\inv(\zeta)}\hat{E}^+_{\bar{\zeta}(\ii)\bar{\tau}(\jj)} \otimes e_{\tau(\jj)\sigma_{\ii}(\tau)(\ii)}.
\]
By definition, $\tau = \sigma_{\ii}(\tau) \xi_{\ii}(\tau)$ and $d_{\ii}(\tau) = \inv(\xi_{\ii}(\tau))$, so therefore
\[
R_{0,\{1,\dots,m\}} = \sum_{\ii,\jj \in \mathcal{W}_m } \sum_{\zeta \in D_{\ii},\tau \in D_{\jj}}q^{\inv(\tau) -d_{\ii}(\tau) -\inv(\zeta)}\hat{E}^+_{\bar{\zeta}(\ii)\bar{\tau}(\jj)} \otimes e_{\tau(\jj)\tau(\ii)}.
\]
By Lemma \ref{prop 1} and the identity $\inv(\bar{\tau}) = m(m-1)/2 - \inv(\tau)$,
\[
\sum_{\zeta \in D_{\ii}}q^{\inv(\tau) -\inv(\zeta)}\hat{E}_{\bar{\zeta}(\ii)\bar{\tau}(\jj)}
\]
does not depend on $\tau$. Therefore, the $R$-matrix factors as
\begin{align*}
R_{0,\{1,\dots,m\} } &= \sum_{\ii,\jj \in \mathcal{W}_m } \bigg( \sum_{\zeta \in D_{\ii}}q^{\inv(\tau) -\inv(\zeta)}\hat{E}^+_{\bar{\zeta}(\ii)\bar{\tau}(\jj)} \bigg) \otimes \bigg( \sum_{\tau \in D_{\jj}} q^{-d_{\ii}(\tau)}e_{\tau(\jj)\tau(\ii)} \bigg) \\
&= \sum_{\ii,\jj \in \mathcal{W}_m } \tilde{E}_{\ii\jj}^+ \otimes e_{\jj\ii}.
\end{align*}
A similar argument shows that
\[
R^T_{0,\{1,\dots,m\}}=\sum_{\ii,\jj \in \mathcal{W}_m } \tilde{E}^-_{\jj\ii} \otimes e_{\ii\jj}.
\]

Therefore,
\[
\Gamma_m = R^T_{0,\{1,\dots,m\} } R_{0,\{1,\dots,m\} } = \sum_{\ii,\jj \in \mathcal{W}_m } \sum_{\ii',\jj' \in \mathcal{W}_m } \tilde{E}_{\jj \ii}^-\tilde{E}_{\ii'\jj'}^+ \otimes e_{\ii\jj}e_{\jj'\ii'}.
\]
By Corollary \ref{Core},
\[
C_m = (\id \otimes \tr_q)(\Gamma_m) = \sum_{\ii,\jj \in \mathcal{W}_m } \tilde{E}_{\jj \ii}^-\tilde{E}_{\ii\jj}^+ \otimes \tr_q(e_{\ii\jj}e_{\jj\ii}).
\]

It just remains to calculate the quantum trace. It is given by
\[
\tr_q(e_{\ii\jj}e_{\jj\ii}) = \tr_q(e_{\ii\ii}) = \tr\big(q^{-NE_{00} -(N-2)E_{11} + \dots + NE_{NN}} e_{\ii\ii}\big) = q^{-N\mu_{0} -(N-2)\mu_{1} + \dots + N\mu_{N}},
\]
where $\mu=\mu^{(N)}(\ii)$. Applying Lemma \ref{Leem} finishes the proof.

\begin{thebibliography}{99}
\bibitem[1]{Arnaudon:1993yg}
Arnaudon D., Bauer M., Polynomial relations in the centre of
 {$\mathcal{U}_q(\mathfrak{sl}(N))$}, \href{https://doi.org/10.1007/BF00805857}{\textit{Lett. Math. Phys.}} \textbf{30}
 (1994), 251--257, \href{https://arxiv.org/abs/hep-th/9310030}{arXiv:hep-th/9310030}.

\bibitem[4]{Carter}
Carter R.W., Finite groups of {L}ie type: conjugacy classes and complex
 characters, Pure Appl. Math. (New York), John Wiley \& Sons, Inc., New York,
 1985.

\bibitem[6]{Drin90}
Drinfeld V.G., Almost cocommutative {H}opf algebras, \textit{Leningrad
 Math.~J.} \textbf{1} (1990), 321--342.

\bibitem[7]{FaddeevRT}
Faddeev L.D., Reshetikhin N.Yu., Takhtajan L.A., Quantization of {L}ie groups
 and {L}ie algebras, in Algebraic Analysis, {V}ol.~{I}, Vol.~1, \href{https://doi.org/10.1016/B978-0-12-400465-8.50019-5}{Academic
 Press}, Boston, MA, 1988, 129--139.

\bibitem[8]{GZB}
Gould M.D., Zhang R.B., Bracken A.J., Generalized {G}el'fand invariants and
 characteristic identities for quantum groups, \href{https://doi.org/10.1063/1.529152}{\textit{J.~Math. Phys.}}
 \textbf{32} (1991), 2298--2303.

\bibitem[9]{Jimbo1986}
Jimbo M., A {$q$}-analogue of {$U(\mathfrak{gl}(N+1))$}, {H}ecke algebra, and
 the {Y}ang--{B}axter equation, \href{https://doi.org/10.1007/BF00400222}{\textit{Lett. Math. Phys.}} \textbf{11} (1986),
 247--252.

\bibitem[10]{KT91}
Khoroshkin S.M., Tolstoy V.N., Universal {$R$}-matrix for quantized
 (super)algebras, \href{https://doi.org/10.1007/BF02102819}{\textit{Comm. Math. Phys.}} \textbf{141} (1991), 599--617.

\bibitem[11]{KirillovResh}
Kirillov A.N., Reshetikhin N., {$q$}-{W}eyl group and a multiplicative formula
 for universal {$R$}-matrices, \href{https://doi.org/10.1007/BF02097710}{\textit{Comm. Math. Phys.}} \textbf{134} (1990),
 421--431.

\bibitem[14]{KIMRN}
Kuan J., A multi-species {${\rm ASEP}(q,j)$} and {$q$}-{TAZRP} with stochastic
 duality, \href{https://doi.org/10.1093/imrn/rnx034}{\textit{Int. Math. Res. Not.}} \textbf{2018} (2018), 5378--5416,
 \href{https://arxiv.org/abs/1605.00691}{arXiv:1605.00691}.

\bibitem[15]{KMMO}
Kuniba A., Mangazeev V.V., Maruyama S., Okado M., Stochastic {$R$} matrix for
 {$U_q\big(A_n^{(1)}\big)$}, \href{https://doi.org/10.1016/j.nuclphysb.2016.09.016}{\textit{Nuclear Phys.~B}} \textbf{913} (2016),
 248--277, \href{https://arxiv.org/abs/1604.08304}{arXiv:1604.08304}.

\bibitem[18]{Lusztig90}
Lusztig G., Canonical bases arising from quantized enveloping algebras,
 \href{https://doi.org/10.2307/1990961}{\textit{J.~Amer. Math. Soc.}} \textbf{3} (1990), 447--498.

\bibitem[21]{Tanisaki1992KILLINGFH}
Tanisaki T., Killing forms, {H}arish-{C}handra isomorphisms, and universal
 {$R$}-matrices for quantum algebras, \href{https://doi.org/10.1142/s0217751x92004117}{\textit{Internat.~J. Modern Phys.~A}}
 \textbf{7} (1992), supp.~01b, 941--961.

\bibitem[22]{Xi1994}
Xi N.H., Root vectors in quantum groups, \href{https://doi.org/10.1007/BF02564506}{\textit{Comment. Math. Helv.}}
 \textbf{69} (1994), 612--639.

\bibitem[23]{ZGBJPhysA}
Zhang R.B., Gould M.D., Bracken A.J., Generalized {G}el'fand invariants of
 quantum groups, \href{https://dx.doi.org/10.1088/0305-4470/24/5/009}{\textit{J.~Phys.~A}} \textbf{24} (1991), 937--943.

\bibitem[24]{ZGBCMP}
Zhang R.B., Gould M.D., Bracken A.J., Quantum group invariants and link
 polynomials, \href{https://doi.org/10.1007/BF02099115}{\textit{Comm. Math. Phys.}} \textbf{137} (1991), 13--27.

\end{thebibliography}
\end{document}
